\documentclass{ltjsarticle}

\usepackage{amsmath}
\usepackage{amssymb}
\usepackage{amsthm}
\usepackage{ mathrsfs }
\usepackage{enumerate}
\usepackage{bbm}
\usepackage{lipsum}
\usepackage{fancyhdr}
\usepackage{tikz-cd} 
\usetikzlibrary{graphs}
\usepackage{luatexja-ruby} 
\usepackage{graphicx} 
\graphicspath{ {./images/} }

\newtheorem{theorem}{定理}[section] 
\newtheorem{proposition}{命題}[section] 
\newtheorem{definition}{定義}[section] 
\newtheorem{problem}{問題}
\newtheorem{postulate}{公準}[section] 
\newtheorem{lemma}{Lemma}[section] 
\newtheorem{notation}{Notation}[section] 
\newtheorem{remark}{注釈}[section] 
\newtheorem{corollary}{系}[section] 
\newtheorem{terminology}{Terminology}[section] 
\newtheorem{example}{例}[section] 
\newtheorem{step}{手順} 
\newtheorem{solution}{解答}[section] 
\newtheorem{fact}{事実}[section] 


\DeclareMathOperator{\diam}{diam}
\DeclareMathOperator{\Gal}{Gal}
\DeclareMathOperator{\rank}{rank}
\DeclareMathOperator{\Hom}{Hom}
\DeclareMathOperator{\Dom}{Dom}
\DeclareMathOperator{\Ann}{Ann}
\DeclareMathOperator{\grad}{grad}
\DeclareMathOperator{\Span}{Span}
\DeclareMathOperator{\interior}{int}
\DeclareMathOperator{\ind}{ind}
\DeclareMathOperator{\supp}{supp}
\DeclareMathOperator{\ob}{ob}
\DeclareMathOperator{\Spec}{Spec}
\DeclareMathOperator{\MaxSpec}{MaxSpec}
\DeclareMathOperator{\PreSh}{PreSh}
\DeclareMathOperator{\Sh}{Sh}
\DeclareMathOperator{\Fun}{Fun}
\DeclareMathOperator{\Ext}{Ext}
\DeclareMathOperator{\Aut}{Aut}
\DeclareMathOperator{\Ker}{Ker}
\DeclareMathOperator{\Image}{Im}
\DeclareMathOperator{\Coker}{Coker}
\DeclareMathOperator{\id}{id}
\DeclareMathOperator{\Mod}{Mod}
\DeclareMathOperator{\QCoh}{QCoh}
\DeclareMathOperator{\Coh}{Coh}
%\newcommand*{\name}[\num_arguments][default values]{{\color{#1}\Large #2}}
\newcommand*{\Z}{\mathbb{Z}}
\newcommand*{\Q}{\mathbb{Q}}
\newcommand*{\R}{\mathbb{R}}
\newcommand*{\multivar}[2]{{#1_1,\cdots,#1_{#2}}}
\newcommand*{\localization}[2]{#1_{\mathfrak{#2}}}
\newcommand*{\ringedspacemorph}[1]{(#1,#1^{\#})}
\newcommand*{\sheaf}[2]{(#1,{\mathcal{#2}}_{#1})}
\newcommand{\fib}[1]{%
  \mathbin{\mathop{\times}\limits_{#1}}%
}
\newcommand{\tens}[1]{%
  \mathbin{\mathop{\otimes}\displaylimits_{#1}}%
}
\title{整数論}
\author{ボン補習校高校数学}
\date{}

\begin{document}

\maketitle

\section{合同式}
\subsection{余のある計算}

\begin{definition}
$n$を自然数とする。整数$a,b$が$n$を法として合同とは
\begin{equation*}
b-a=nk
\end{equation*}
となる整数$k$が存在することである。これを
\begin{equation*}
a\equiv b \mod{n}
\end{equation*}
と書く。
\end{definition}

次の三つは全て同じ意味です。
\begin{center}
\begin{enumerate}[1).]
\item $a,b$は$n$を法として合同。
\item $b-a$は$n$の倍数。
\item $a,b$を$n$で割った時の余りが同じ。
\end{enumerate}
\end{center}

このように、別の言葉を使って数学的に全く同じことを言える時、それらの条件は\ltjruby{同値}{どうち}であると言います。

\begin{example}他の分野でも上のような同値な条件があります。たとえば、
\begin{enumerate}[1).]
\item 三角形$ABC$は正三角形である（それぞれの辺の長さが等しい）。
\item 三角形$ABC$の三つの角度が全て等しい（それぞれ$60^\circ$である）。
\end{enumerate}
\end{example}

\begin{definition}
整数$p$が素数であるとは、どんな二つの整数$a,b$に対して$ab$が$p$の倍数ならば$a$もしくは$b$が$p$の倍数である時のことである。
\end{definition}

上の定義の例を見てみましょう。例えば$p=7$の時、$a=10,b=9$としたら$ab=90$は$7$の倍数でないので$a,b$どちらも$7$の倍数ではありません。逆に$2\times 3=6$は素数ではないです。$a=4,b=9$とすると$ab=36$で$6$の倍数になりますが$4,9$どちらも$6$の倍数ではありません。これを合同式を使って言い換えると次のようになります。
\begin{proposition}
$p$が素数であるとは
\begin{equation*}
ab\equiv 0\mod{p}
\end{equation*}
ならば
\begin{equation*}
a\equiv 0\mod{p}\text{もしくは}b\equiv0\mod{p}
\end{equation*}
となる。
\label{prop_prime_element_mod}
\end{proposition}

\begin{proposition}
$p$を素数とする。整数$a,b,c$について考える。もし
\begin{equation*}
a\not\equiv 0\mod{p}
\end{equation*}
で
\begin{equation*}
ab \equiv ac\mod{p}
\end{equation*}
ならば
\begin{equation*}
b\equiv c\mod{p}
\end{equation*}
\label{prop_injectivity_multiplication}
\end{proposition}

\begin{proof}[証明]条件を使って
\begin{equation*}
ab-ac\equiv0\mod{p}
\end{equation*}
となる。よって
\begin{equation*}
a(b-c)\equiv 0 \mod{p}
\end{equation*}
ここで$p$は素数かつ$a\equiv0\mod{p}$より
\begin{equation*}
b-c\equiv0\mod{p}\Rightarrow b\equiv c\mod{p}
\end{equation*}
\end{proof}

$p=5$の時を考えます。あまりの可能性は$0,1,2,3,4$です。これらにそれぞれ$3$をかけてみます。
\begin{align*}
3\times 0 &\equiv 0 \mod{5}\\
3\times 1 &\equiv 3 \mod{5}\\
3\times 2 &\equiv 6\equiv 1 \mod{5}\\
3\times 3 &\equiv 9\equiv4 \mod{5}\\
3\times 4 &\equiv 12\equiv 2 \mod{5}\\
\end{align*}

このように五つそれぞれ異なる余りになりました。命題\ref{prop_injectivity_multiplication}の言っていることが少しわかりやすくなったと思います。ここで三番目の
\begin{equation*}
3\times 2 \equiv 6\equiv 1 \mod{5}\\
\end{equation*}

に注目してみます。このように掛け算をしたら$1$となるペアはどのようなあまりにも存在するでしょうか。
\begin{proposition}
$p$を素数とする。もし$a\equiv0\mod{p}$なら$ab\equiv1\mod{1}$となる整数$b$がある。
\end{proposition}
\begin{proof}[証明]
命題\ref{prop_injectivity_multiplication}を使う。$p=5$の時に見たように
\begin{equation*}
a\times 0,a\times 1,\cdots,a\times (p-1)
\end{equation*}
はそれぞれ$p$個の異なる余りを持つ。あまりの可能性も$p$個しかないので$a\times b\equiv1\mod{p}$となる$b$がある。
\end{proof}
同じ論理で次のことも言えます。
\begin{proposition}
$p$を素数とする。$a\equiv0\mod{p}$なら$ab\equiv1\mod{p}$となる$b$は一つしかない。
\end{proposition}
\begin{proposition}
$p$を素数とする。$p$を法として自分自身が逆数となる数は$1,-1$のみである。
\end{proposition}
\begin{proof}[証明]
自分自身が逆数になるとは
\begin{equation*}
x^2=1\mod{p}
\end{equation*}
となることである。$1$を左辺に移して
\begin{equation*}
x^2-1\equiv 0\mod{p}
\end{equation*}
左辺を因数分解して
\begin{equation*}
(x-1)(x+1)\equiv 0\mod{p}
\end{equation*}
命題\ref{prop_prime_element_mod}を使うと
\begin{equation*}
x-1\equiv0\mod{p}\text{もしくは}x+1\equiv 0\mod{p}
\end{equation*}
よって$x\equiv\pm1\mod{p}$。
\end{proof}
\begin{theorem}[Wilsonの定理]
$p$を素数とする。次のことが成り立つ。
\begin{equation*}
(p-1)!\equiv -1\mod{p}
\end{equation*}
\end{theorem}
\begin{theorem}[Fermatの小定理]
$p$を素数とする。次のことが成り立つ。
\begin{equation*}
a^{p-1}\equiv 1\mod{p}
\end{equation*}
\end{theorem}
\begin{proof}

\end{proof}
\end{document}